\documentclass[12pt]{article}
\usepackage[utf8]{inputenc}
\usepackage[russian]{babel}
\usepackage{amsmath, amssymb, amsthm}
\usepackage{graphicx}
\usepackage{hyperref}
\usepackage{geometry}
\geometry{a4paper, margin=2.5cm}

\newtheorem{definition}{Определение}
\newtheorem{theorem}{Теорема}
\newtheorem{lemma}{Лемма}
\newtheorem{corollary}{Следствие}
\newtheorem{example}{Пример}
\newtheorem{algorithm}{Алгоритм}

\title{Конспект по теории графов и бинарным отношениям}
\author{По материалам учебных пособий}
\date{}

\begin{document}

\maketitle
\tableofcontents
\newpage

\section{Бинарные отношения и графы}
\subsection{Основные понятия}
Пусть \(A\) — произвольное множество. \textbf{Бинарным отношением} \(P\) на множестве \(A\) называется любое подмножество декартова произведения \(A \times A\), т.е. \(P \subseteq A \times A\). Если пара \((x,y) \in P\), то пишут \(xPy\) и говорят, что \(x\) находится в отношении \(P\) с \(y\).

\subsection{Способы задания бинарных отношений}
\begin{enumerate}
    \item \textbf{Перечислением:} \(P = \{(x_1,y_1), (x_2,y_2), \dots\}\).
    \item \textbf{Матрицей смежности:} квадратная матрица \(\|a_{xy}\|\) размера \(|A| \times |A|\), где \(a_{xy} = 1\) если \((x,y) \in P\), и \(0\) иначе.
    \item \textbf{Графом:} ориентированный граф \(G = (A, P)\), в котором вершины — элементы \(A\), а дуга \((x,y)\) существует тогда и только тогда, когда \((x,y) \in P\).
\end{itemize}

\subsection{Свойства бинарных отношений}
Бинарное отношение \(P\) на множестве \(A\) называется:
\begin{itemize}
    \item \textbf{рефлексивным}, если \(\forall x\in A: (x,x)\in P\);
    \item \textbf{антирефлексивным}, если \(\forall x\in A: (x,x)\notin P\);
    \item \textbf{симметричным}, если \(\forall x,y\in A: xPy \Rightarrow yPx\);
    \item \textbf{асимметричным}, если \(\forall x,y\in A: xPy \Rightarrow \neg (yPx)\);
    \item \textbf{транзитивным}, если \(\forall x,y,z\in A: xPy \land yPz \Rightarrow xPz\);
    \item \textbf{отрицательно транзитивным}, если \(\forall x,y,z\in A: \neg(xPy) \land \neg(yPz) \Rightarrow \neg(xPz)\);
    \item \textbf{связным}, если \(\forall x\neq y: xPy \text{ или } yPx\);
    \item \textbf{полным}, если \(\forall x,y: xPy \text{ или } yPx\).
\end{itemize}

\subsection{Специальные классы бинарных отношений}
\begin{itemize}
    \item \textbf{Частичный порядок} — антирефлексивное, транзитивное и ациклическое отношение.
    \item \textbf{Слабый порядок} — частичный порядок, удовлетворяющий условию отрицательной транзитивности.
    \item \textbf{Линейный порядок} — связный слабый порядок.
\end{itemize}

\subsection{Теорема о представлении слабого порядка функцией полезности}
\begin{theorem}[Представление слабого порядка]
    Пусть \(A\) — конечное множество. Бинарное отношение \(P\) на \(A\) является слабым порядком тогда и только тогда, когда существует функция \(u: A \to \mathbb{R}\) такая, что
    \[
    xPy \iff u(x) > u(y).
    \]
\end{theorem}
\begin{proof}[Доказательство (набросок)]
    \textit{Необходимость:} Пусть \(P\) — слабый порядок. Определим \(u(x) = |\{y \in A : xPy\}|\) (количество элементов, которые хуже \(x\)). Тогда \(xPy\) влечёт \(u(x) > u(y)\) благодаря транзитивности и антирефлексивности. Обратно, если \(u(x) > u(y)\), то из условия Чипмана (эквивалентного отрицательной транзитивности) следует \(xPy\).
    
    \textit{Достаточность:} Если \(xPy \iff u(x) > u(y)\), то \(P\) антирефлексивно, транзитивно и отрицательно транзитивно, т.е. является слабым порядком.
\end{proof}

\section{Способы представления графов}
Граф \(G = (V, E)\) может быть задан:
\begin{enumerate}
    \item \textbf{Матрицей смежности} — квадратная матрица \(n \times n\) (\(n = |V|\)), где элемент \(a_{ij} = 1\), если \((v_i, v_j) \in E\), и \(0\) иначе. Для неориентированных графов матрица симметрична.
    \item \textbf{Списком рёбер} — перечисление всех пар вершин, соединённых ребром.
    \item \textbf{Списком смежности} — для каждой вершины \(v\) указывается список всех вершин, смежных с \(v\).
\end{enumerate}

\section{Подграфы. Маршруты, цепи, циклы. Связность}
\subsection{Подграф}
Подграфом графа \(G = (V,E)\) называется граф \(G' = (V', E')\), где \(V' \subseteq V\), \(E' \subseteq E\) и все концы рёбер из \(E'\) лежат в \(V'\).

\subsection{Маршруты, цепи, циклы}
\begin{itemize}
    \item \textbf{Маршрут} — последовательность вершин \(v_1, v_2, \dots, v_k\) такая, что каждые две соседние соединены ребром (в ориентированном графе — дугой).
    \item \textbf{Цепь} — маршрут, в котором все рёбра различны.
    \item \textbf{Простая цепь (путь)} — цепь, в которой все вершины различны.
    \item \textbf{Цикл} — замкнутая цепь (\(v_1 = v_k\)), в которой все рёбра различны.
    \item \textbf{Простой цикл} — цикл, в котором все вершины различны (кроме первой и последней).
\end{itemize}

\subsection{Связность}
\begin{itemize}
    \item \textbf{Неориентированный граф} называется \textbf{связным}, если для любых двух вершин существует цепь, соединяющая их.
    \item \textbf{Компонента связности} — максимальный связный подграф.
    \item \textbf{Ориентированный граф} называется \textbf{сильно связным}, если для любых двух вершин \(u,v\) существует ориентированный путь из \(u\) в \(v\) и из \(v\) в \(u\). \textbf{Слабо связным} — если его неориентированный аналог связен.
\end{itemize}

\section{Эйлеровы пути и циклы}
\subsection{Определения}
\begin{itemize}
    \item \textbf{Эйлеров цикл} — цикл, проходящий по каждому ребру графа ровно один раз.
    \item \textbf{Эйлеров путь} — цепь, проходящая по каждому ребру ровно один раз.
    \item Граф, обладающий эйлеровым циклом, называется \textbf{эйлеровым}.
\end{itemize}

\subsection{Теорема Эйлера}
\begin{theorem}[Эйлер, 1736]
    Связный неориентированный граф обладает эйлеровым циклом тогда и только тогда, когда степени всех его вершин чётны.
    Связный неориентированный граф обладает эйлеровым путём (не циклом) тогда и только тогда, когда ровно две вершины имеют нечётную степень.
\end{theorem}
\begin{proof}
    \textit{Необходимость:} При обходе цикла каждое вхождение в вершину сопровождается выходом, поэтому степень чётна. Для пути нечётными будут начальная и конечная вершины.
    
    \textit{Достаточность:} (Конструктивное доказательство) Начинаем из любой вершины и идём по рёбрам, пока не вернёмся в начальную. Получим цикл \(C\). Если остались непройденные рёбра, найдём вершину \(v\) на \(C\), инцидентную таким рёбрам. Построим новый цикл \(C'\), начинающийся в \(v\) и использующий только непройденные рёбра (чётность степеней гарантирует возврат). Объединим \(C\) и \(C'\) в один цикл. Повторяем, пока все рёбра не будут включены.
\end{proof}

\subsection{Алгоритм Флёри (Fleury)}
\begin{algorithm}[Алгоритм Флёри]
    \textbf{Идея:} избегать использования \textit{мостов} (рёбер, удаление которых увеличивает число компонент связности), если только нет другого выбора.
    \begin{enumerate}
        \item Начать с произвольной вершины (если нужен эйлеров путь — с одной из нечётных вершин).
        \item На каждом шаге выбирать ребро, удаление которого не делает граф несвязным (т.е. не мост), если только нет альтернативы. Если все оставшиеся рёбра — мосты, то разрешается идти по мосту.
        \item Пройти по выбранному ребру, удалить его.
        \item Повторять, пока есть рёбра.
    \end{enumerate}
    \textbf{Почему работает:} Мосты могут быть пройдены только в конце, чтобы не разбить граф на компоненты. Алгоритм корректен, но требует проверки на мост на каждом шаге, что можно сделать за \(O(|E|)\) (например, поиском в глубину), поэтому общая сложность \(O(|E|^2)\).
\end{algorithm}

\subsection{Алгоритм Хирхольцера (Hierholzer)}
\begin{algorithm}[Алгоритм Хирхольцера, 1873]
    Более эффективный, чем Флёри, имеет линейную сложность \(O(|E|)\).
    \begin{enumerate}
        \item Выбрать стартовую вершину \(v_0\) (если нужен эйлеров путь — одну из нечётных вершин).
        \item Построить цикл \(C\), начиная с \(v_0\) и идя по непройденным рёбрам, пока не вернёмся в \(v_0\) (это всегда возможно благодаря чётности степеней).
        \item Если в \(C\) есть вершина \(u\), у которой остались непройденные рёбра, начать новый цикл \(C'\) из \(u\), используя только непройденные рёбра.
        \item Вставить \(C'\) в \(C\) на место \(u\).
        \item Повторять, пока все рёбра не будут включены.
    \end{enumerate}
    Реализация использует стек и обход в глубину. Сложность \(O(|E|)\).
\end{algorithm}

\subsection{Эйлеровы пути в ориентированных графах}
\begin{theorem}
    Связный ориентированный граф (слабо связный) имеет эйлеров цикл тогда и только тогда, когда для каждой вершины \(v\) полустепень захода равна полустепени исхода: \(\deg^+(v) = \deg^-(v)\). Эйлеров путь существует, если ровно одна вершина имеет \(\deg^+ = \deg^- + 1\) (начало), ровно одна — \(\deg^- = \deg^+ + 1\) (конец), а у остальных равенство.
\end{theorem}
\begin{algorithm}[Алгоритм Хирхольцера для ориентированных графов]
    Аналогичен неориентированному случаю, но обход всегда идёт по направлению дуг. Сложность \(O(|E|)\).
\end{algorithm}

\section{Гамильтоновы пути и циклы}
\subsection{Определения}
\begin{itemize}
    \item \textbf{Гамильтонов цикл} — простой цикл, проходящий через каждую вершину графа ровно один раз.
    \item \textbf{Гамильтонов путь} — простая цепь, проходящая через все вершины.
\end{itemize}

\subsection{Сложность задачи}
Задача проверки существования гамильтонова цикла в общем графе является \textbf{NP-полной} (Карп, 1972). 
\begin{itemize}
    \item \textbf{Что такое NP-полнота?} Класс NP (nondeterministic polynomial time) — задачи, решение которых можно \textit{проверить} за полиномиальное время, если предъявить сертификат (например, сам гамильтонов цикл). NP-полные задачи — самые «трудные» в NP: если для какой-то NP-полной задачи найден полиномиальный алгоритм, то и все задачи NP решаются полиномиально (P = NP). 
    \item \textbf{Почему гамильтонов цикл NP-полный?} Потому что к нему можно свести любую задачу из NP за полиномиальное время (доказательство Карпа). Для проверки гамильтонова цикла достаточно предъявить порядок обхода вершин и убедиться, что каждые две соседние соединены ребром — это делается за \(O(n)\).
    \item \textbf{Что нужно для проверки?} Перебор всех \(n!\) перестановок вершин — экспоненциальный алгоритм. Полиномиальный алгоритм для произвольных графов неизвестен.
\end{itemize}
Для некоторых классов (полные графы, турниры, планарные кубические графы) существуют полиномиальные алгоритмы (например, для турниров — всегда есть гамильтонов путь).

\section{Раскраска графов}
\subsection{Основные понятия}
\begin{itemize}
    \item \textbf{Раскраска вершин} — приписывание вершинам цветов так, что любые две смежные вершины имеют разные цвета.
    \item \textbf{Хроматическое число} \(\chi(G)\) — минимальное количество цветов, необходимое для раскраски графа.
    \item \textbf{Планарный граф} — граф, который можно нарисовать на плоскости без пересечения рёбер (кроме вершин).
    \item \(\delta(G)\) — минимальная степень вершины в графе \(G\).
\end{itemize}

\subsection{Теорема о пяти красках (для плоских графов)}
\begin{theorem}[Хивуд, 1890]
    Всякий планарный граф (карту) можно раскрасить в 5 цветов.
\end{theorem}
\begin{proof}[Идея доказательства]
    Индукция по числу вершин. Для планарного графа всегда существует вершина степени \(\le 5\) (следует из формулы Эйлера). Удаляем её, раскрашиваем остаток в 5 цветов (по предположению индукции). Возвращаем вершину. Если у неё менее 5 соседей, цвет найдётся. Если ровно 5 соседей, то рассматриваем подграф, индуцированный двумя цветами, и перекрашиваем.
\end{proof}

\subsection{Проблема четырёх красок}
Любой планарный граф можно раскрасить в 4 цвета. Доказана с помощью компьютера (Аппель, Хакен, 1976).

\section{Двудольные графы. Паросочетания}
\subsection{Двудольные графы}
Граф \(G = (V,E)\) называется \textbf{двудольным}, если множество вершин можно разбить на две доли \(X\) и \(Y\) так, что все рёбра соединяют вершины из разных долей.

\begin{theorem}[Критерий двудольности]
    Граф является двудольным тогда и только тогда, когда он не содержит циклов нечётной длины.
\end{theorem}

\subsection{Паросочетания}
\textbf{Паросочетание} \(M\) — множество рёбер, не имеющих общих вершин. \textbf{Максимальное паросочетание} — наибольшее по числу рёбер. \textbf{Совершенное паросочетание} — покрывает все вершины.

\begin{example}[Совершенное паросочетание]
    В полном двудольном графе \(K_{2,2}\) (вершины \(X=\{x_1,x_2\}\), \(Y=\{y_1,y_2\}\)) паросочетание \(\{(x_1,y_1), (x_2,y_2)\}\) является совершенным. Оно может существовать только если \(|X|=|Y|\). В связном неориентированном графе совершенное паросочетание возможно, например, в цикле \(C_4\) (квадрат) — он связен, и паросочетание из двух противоположных рёбер совершенно.
\end{example}

\subsection{Теорема Холла}
\begin{theorem}[Холл, 1935]
    В двудольном графе \(G = (X \cup Y, E)\) существует паросочетание, покрывающее все вершины \(X\), тогда и только тогда, когда для любого подмножества \(S \subseteq X\) выполняется
    \[
    |N(S)| \ge |S|,
    \]
    где \(N(S)\) — множество соседей \(S\) в \(Y\).
\end{theorem}

\subsection{Алгоритм построения наибольшего паросочетания (алгоритм Куна)}
\begin{algorithm}
    \begin{enumerate}
        \item Начать с пустого паросочетания.
        \item Для каждой вершины \(x \in X\) пытаться найти \textit{чередующуюся цепь} (начинающуюся в свободной вершине \(x\) и заканчивающуюся свободной вершиной в \(Y\)).
        \item Если такая цепь найдена, увеличить паросочетание, меняя рёбра вдоль цепи.
        \item Повторять, пока возможно.
    \end{enumerate}
    Сложность \(O(|V| \cdot |E|)\).
\end{algorithm}

\section{Деревья}
\subsection{Эквивалентные определения}
Дерево — связный граф без циклов. Эквивалентные условия:
\begin{enumerate}
    \item Граф связен и не содержит циклов.
    \item Граф связен и имеет \(|V| - 1\) ребро.
    \item Граф ацикличен и имеет \(|V| - 1\) ребро.
    \item Между любыми двумя вершинами существует единственный простой путь.
\end{enumerate}

\begin{theorem}[Уникальность пути]
    В дереве между любыми двумя вершинами существует ровно один простой путь.
\end{theorem}
\begin{proof}
    От противного: если бы было два пути, то они образовывали бы цикл, что противоречит определению дерева.
\end{proof}

\subsection{Способы представления деревьев}
\begin{itemize}
    \item \textbf{Корневое дерево} — выделена вершина-корень, все рёбра ориентированы от корня.
    \item \textbf{Двоичное дерево} — у каждой вершины не более двух детей.
    \item \textbf{Сбалансированное (AVL) дерево} — для каждой вершины высоты поддеревьев различаются не более чем на 1.
\end{itemize}

\begin{example}[Сбалансированные деревья]
    \begin{itemize}
        \item \textbf{AVL-дерево}: после вставки/удаления выполняется балансировка поворотами. Высота \(O(\log n)\).
        \item \textbf{Красно-чёрное дерево}: каждый узел красный или чёрный, корень чёрный, красные узлы не могут быть рядом, все пути от корня до листьев содержат одинаковое число чёрных узлов.
        \item \textbf{Двоичная куча} (heap) — не обязательно сбалансирована по высоте, но является полным деревом.
    \end{itemize}
    \textbf{Построение}: для AVL — вставка с последующими поворотами (LL, RR, LR, RL). Для красно-чёрного — вставка с перекрашиванием и вращениями.
\end{example}

\section{Алгоритмы на графах}
\subsection{Обходы графа}
\begin{itemize}
    \item \textbf{Поиск в глубину (DFS)} — рекурсивный обход: посещаем вершину, затем рекурсивно всех её непосещённых соседей. Время \(O(|V| + |E|)\).
    \item \textbf{Поиск в ширину (BFS)} — обход по уровням: посещаем все вершины на расстоянии \(k\) перед \(k+1\). Использует очередь. Время \(O(|V| + |E|)\).
\end{itemize}

\subsection{Топологическая сортировка}
\begin{definition}
    \textbf{Топологическая сортировка} ориентированного ациклического графа (DAG) — линейное упорядочение вершин, в котором для каждого ребра \((u,v)\) вершина \(u\) идёт раньше \(v\).
\end{definition}
\begin{algorithm}[Алгоритм Кана (Kahn)]
    \begin{enumerate}
        \item Вычислить полустепени захода для всех вершин.
        \item Поместить все вершины с нулевой полустепенью захода в очередь.
        \item Пока очередь не пуста:
            \begin{itemize}
                \item Извлечь вершину \(u\), добавить её в ответ.
                \item Для каждого соседа \(v\) из \(u\) уменьшить \(\deg^-(v)\) на 1; если стало 0, добавить \(v\) в очередь.
            \end{itemize}
        \item Если в ответе не все вершины — граф содержит цикл.
    \end{enumerate}
    Сложность \(O(|V|+|E|)\).
\end{algorithm}
Применения: планирование задач, вычисление порядка компиляции.

\subsection{Поиск сильно связных компонент}
\begin{definition}
    \textbf{Сильно связная компонента} (SCC) — максимальное множество вершин, в котором для любой пары \(u,v\) есть пути из \(u\) в \(v\) и из \(v\) в \(u\).
\end{definition}
\begin{algorithm}[Алгоритм Косарайю (Kosaraju)]
    \begin{enumerate}
        \item Выполнить DFS на исходном графе, записывая порядок окончания обработки вершин (post-order).
        \item Построить обратный граф \(G^T\) (все рёбра наоборот).
        \item Обходить вершины в порядке убывания времени окончания из шага 1, запускать DFS на \(G^T\). Каждое дерево DFS — одна SCC.
    \end{enumerate}
    Сложность \(O(|V|+|E|)\).
\end{algorithm}
\begin{algorithm}[Алгоритм Тарьяна (Tarjan)]
    Один проход DFS, использует индексы и стек. Сложность \(O(|V|+|E|)\).
\end{algorithm}

\subsection{Двусвязность, мосты, точки сочленения}
\begin{definition}
    \begin{itemize}
        \item \textbf{Мост} — ребро, удаление которого увеличивает число компонент связности.
        \item \textbf{Точка сочленения} — вершина, удаление которой увеличивает число компонент связности.
        \item Граф называется \textbf{двусвязным}, если он связен и не содержит точек сочленения.
    \end{itemize}
\end{definition}
\begin{algorithm}[Алгоритм Тарьяна для поиска мостов и точек сочленения]
    Использует DFS, вычисляет для каждой вершины \(tin\) (время входа) и \(low\) (минимальное \(tin\) достижимое через обратные рёбра). 
    \begin{itemize}
        \item Ребро \((u,v)\) — мост, если \(low[v] > tin[u]\).
        \item Вершина \(u\) (не корень) — точка сочленения, если существует ребро \((u,v)\) в дереве DFS, такое что \(low[v] \ge tin[u]\). Для корня — если у него более одного ребёнка в DFS-дереве.
    \end{itemize}
    Сложность \(O(|V|+|E|)\).
\end{algorithm}

\subsection{Система непересекающихся множеств (DSU)}
\begin{definition}
    \textbf{DSU (Disjoint Set Union)} — структура данных, поддерживающая объединение множеств и запрос принадлежности элемента к множеству. Применяется в алгоритме Краскала для проверки, образуют ли добавляемые рёбра цикл.
\end{definition}

\begin{algorithm}[Базовые операции DSU]
    \begin{itemize}
        \item \textbf{makeSet(x)}: создать новое множество, содержащее элемент \(x\).
        \item \textbf{find(x)}: найти представителя (корень) множества, содержащего \(x\).
        \item \textbf{union(x, y)}: объединить множества, содержащие \(x\) и \(y\).
    \end{itemize}
\end{algorithm}

\subsubsection{Реализация и оптимизации}
Используется лес, где каждый элемент хранит ссылку на родителя. Корень множества указывает сам на себя.

\begin{enumerate}
    \item \textbf{Наивная реализация:} find проходит вверх до корня \(O(n)\), union связывает корни \(O(1)\).
    \item \textbf{Сжатие пути (path compression):} при вызове find перенаправляем ссылки всех пройденных элементов прямо на корень. Амортизированная сложность почти константная.
    \item \textbf{Объединение по рангу (union by rank):} храним высоту дерева (ранг), при объединении подвешиваем дерево с меньшим рангом к большему. Ранг растёт только при равных рангах.
\end{enumerate}

\begin{algorithm}[Псевдокод DSU с оптимизациями]
    \begin{verbatim}
    parent = [0..n-1]
    rank = [0..n-1]
    
    function makeSet(x):
        parent[x] = x
        rank[x] = 0
    
    function find(x):
        if parent[x] != x:
            parent[x] = find(parent[x])  # сжатие пути
        return parent[x]
    
    function union(x, y):
        xRoot = find(x)
        yRoot = find(y)
        if xRoot == yRoot: return
        if rank[xRoot] < rank[yRoot]:
            parent[xRoot] = yRoot
        else if rank[xRoot] > rank[yRoot]:
            parent[yRoot] = xRoot
        else:
            parent[yRoot] = xRoot
            rank[xRoot] += 1
    \end{verbatim}
\end{algorithm}

\textbf{Сложность:} при использовании обеих оптимизаций амортизированное время операций \(O(\alpha(n))\), где \(\alpha(n)\) — обратная функция Аккермана, растущая крайне медленно (для практических \(n\) не превышает 5). Таким образом, можно считать операции почти константными.

\subsection{Поиск кратчайших путей (алгоритм Дейкстры)}
\begin{algorithm}[Дейкстра]
    Находит кратчайшие расстояния от одной вершины до всех остальных в графе с \textbf{неотрицательными} весами.
    \begin{enumerate}
        \item Инициализация: \(dist[s] = 0\), \(dist[v] = \infty\) для \(v \neq s\). Множество непосещённых вершин.
        \item Пока есть непосещённые вершины: выбрать вершину \(u\) с минимальным \(dist[u]\) среди непосещённых.
        \item Для каждого соседа \(v\) из \(u\): если \(dist[u] + w(u,v) < dist[v]\), обновить \(dist[v]\).
        \item Пометить \(u\) как посещённую.
    \end{enumerate}
    С использованием приоритетной очереди (кучи) сложность \(O(|E| + |V|\log|V|)\).
\end{algorithm}
\textbf{Визуализация (описание):} Представьте граф из 5 вершин (A,B,C,D,E) с весами. Алгоритм начинает из A, помечает расстояние 0. Затем рассматривает соседей, обновляет расстояния, выбирает вершину с наименьшим расстоянием (например, C), и так далее. На каждом шаге мы «растем» область уже посещённых вершин, зная точные кратчайшие расстояния до них.

\subsection{Построение минимального остовного дерева}
\begin{itemize}
    \item \textbf{Алгоритм Краскала}: сортируем рёбра по весу, добавляем наименьшие, не образующие цикла (используем DSU). \(O(|E|\log|E|)\).
    \item \textbf{Алгоритм Прима}: начинаем с одной вершины, на каждом шаге добавляем наименьшее ребро, соединяющее дерево с новой вершиной. \(O(|E|\log|V|)\) с кучей.
\end{itemize}

\subsection{Алгоритм Краскала (подробно)}
\begin{algorithm}[Краскал, 1956]
    \begin{enumerate}
        \item Отсортировать все рёбра по неубыванию веса.
        \item Инициализировать DSU для каждой вершины (вызвать makeSet для всех вершин).
        \item Для каждого ребра \((u,v)\) в отсортированном порядке:
            \begin{itemize}
                \item Если find(u) != find(v) (вершины в разных компонентах), добавить ребро в остов и выполнить union(u, v).
                \item Иначе пропустить (ребро образует цикл).
            \end{itemize}
        \item Конец, когда добавлено \(|V|-1\) ребро.
    \end{enumerate}
    Сложность: сортировка \(O(|E|\log|E|)\), DSU почти линейная.
\end{algorithm}
\textbf{Визуализация (описание):} Представьте точки на плоскости. Алгоритм сначала соединяет самые близкие пары, затем следующие, но если ребро замыкает цикл (все его вершины уже соединены), оно пропускается. В итоге получается «скелет» минимальной суммарной длины.

\subsection{Алгоритм Прима (подробно)}
\begin{algorithm}[Прим, 1957]
    \begin{enumerate}
        \item Выбрать произвольную стартовую вершину \(s\). Множество \(U = \{s\}\).
        \item Для всех вершин \(v \notin U\) хранить ключ \(key[v]\) — минимальный вес ребра из \(U\) в \(v\).
        \item Использовать приоритетную очередь (мин-кучу) для вершин вне \(U\) с ключами.
        \item Пока \(U \neq V\):
            \begin{itemize}
                \item Извлечь вершину \(u\) с минимальным ключом из очереди, добавить её в \(U\) и добавить соответствующее ребро в остов.
                \item Для каждого соседа \(v\) вершины \(u\), если \(v \notin U\) и \(w(u,v) < key[v]\), обновить \(key[v]\) и перестроить очередь.
            \end{itemize}
    \end{enumerate}
    Сложность: \(O(|E|\log|V|)\).
\end{algorithm}
\textbf{Визуализация (описание):} Начинаем с одной вершины. Затем «растим» дерево, добавляя ближайшую не присоединённую вершину — так же, как в алгоритме Дейкстры, но с весами рёбер, а не суммой расстояний.

\section{Максимальный поток и минимальный разрез}
\subsection{Основные понятия}
Рассмотрим сеть (ориентированный граф с источником \(s\) и стоком \(t\), каждому ребру приписана пропускная способность \(c(u,v) \ge 0\)). \textbf{Поток} \(f\) удовлетворяет:
\begin{itemize}
    \item \(0 \le f(u,v) \le c(u,v)\);
    \item \(\sum_{v} f(u,v) = \sum_{v} f(v,u)\) для всех \(u \neq s,t\) (закон сохранения).
\end{itemize}
\textbf{Величина потока} \(|f| = \sum_{v} f(s,v)\). \textbf{Разрез} \((S,T)\) — разбиение вершин, где \(s \in S\), \(t \in T\). Его пропускная способность \(c(S,T) = \sum_{u \in S, v \in T} c(u,v)\).

\subsection{Теорема Форда–Фалкерсона}
\begin{theorem}[Форд–Фалкерсон, 1956]
    Максимальная величина потока в сети равна минимальной пропускной способности разреза.
\end{theorem}
\begin{proof}[Доказательство (набросок)]
    \begin{itemize}
        \item \textit{Слабая двойственность:} Для любого потока \(f\) и любого разреза \((S,T)\) имеем \(|f| \le c(S,T)\). Действительно, поток из \(S\) в \(T\) не может превысить пропускную способность разреза.
        \item \textit{Сильная двойственность:} Рассмотрим максимальный поток. Построим остаточную сеть \(G_f\). Если в \(G_f\) нет пути из \(s\) в \(t\), то множество вершин, достижимых из \(s\) в \(G_f\), образует разрез \((S,T)\), причём все рёбра из \(S\) в \(T\) в исходной сети насыщены (\(f(u,v)=c(u,v)\)), а из \(T\) в \(S\) — пусты (\(f(v,u)=0\)). Тогда \(|f| = c(S,T)\). Значит, максимальный поток равен минимальному разрезу.
    \end{itemize}
\end{proof}

\subsection{Алгоритм Форда–Фалкерсона (подробно)}
\begin{algorithm}
    \begin{enumerate}
        \item Начать с нулевого потока \(f(u,v)=0\) для всех рёбер.
        \item Построить остаточную сеть \(G_f\), где для каждого ребра \((u,v)\) есть прямое ребро с остаточной пропускной способностью \(c_f(u,v) = c(u,v) - f(u,v)\) и обратное ребро с \(c_f(v,u) = f(u,v)\).
        \item Пока существует увеличивающий путь (путь из \(s\) в \(t\) в \(G_f\)):
            \begin{itemize}
                \item Найти такой путь \(P\) (например, BFS/DFS).
                \item Вычислить \(\delta = \min_{(u,v) \in P} c_f(u,v)\).
                \item Увеличить поток: для каждого ребра \((u,v)\) на \(P\):
                    \begin{itemize}
                        \item если \((u,v)\) — прямое ребро в исходной сети, то \(f(u,v) := f(u,v) + \delta\);
                        \item если \((u,v)\) — обратное ребро (соответствует \(f(v,u)>0\)), то \(f(v,u) := f(v,u) - \delta\).
                    \end{itemize}
                \item Обновить остаточную сеть.
            \end{itemize}
        \item Конец. Полученный поток максимален.
    \end{enumerate}
    \textbf{Сложность:} \(O(|E| \cdot |f_{\max}|)\) в худшем случае, где \(|f_{\max}|\) — величина максимального потока. Это \textbf{псевдополиномиальная} сложность (зависит от значения, а не от размера входных данных в битах). При целых пропускных способностях алгоритм завершится за конечное число шагов. Если пропускные способности иррациональны, может не сойтись. Для практического использования применяют улучшенные версии (Эдмондс–Карп с BFS, сложность \(O(|V|\cdot|E|^2)\) или Dinic с \(O(|E|\sqrt{|V|})\) для единичных пропускных способностей).
\end{algorithm}

\subsection{Применения потоков}
\begin{itemize}
    \item \textbf{Максимальное паросочетание в двудольном графе} сводится к потоку: из источника в левую долю, из левой в правую (рёбра графа), из правой в сток, все пропускные способности 1. Максимальный поток равен размеру максимального паросочетания.
    \item \textbf{Задача о назначениях} (минимизация стоимости) — транспортная задача, решается алгоритмом минимального потока (или венгерским алгоритмом).
    \item \textbf{Теорема Менгера}: Для любых двух вершин \(u,v\) размер минимального \(u,v\)-разреза (минимальное число рёбер/вершин, удаление которых разделяет \(u\) и \(v\)) равен максимальному числу непересекающихся путей (по рёбрам или по вершинам) между \(u\) и \(v\). Следствие: связность графа равна максимальному числу непересекающихся путей.
\end{itemize}

\section{Введение в теорию сложности}
\subsection{Классы P и NP}
\begin{itemize}
    \item \textbf{Класс P} — задачи, решаемые детерминированной машиной Тьюринга за полиномиальное время от длины входа. Пример: поиск кратчайшего пути.
    \item \textbf{Класс NP} — задачи, решение которых можно \textit{проверить} за полиномиальное время, если предъявить сертификат (доказательство). Пример: гамильтонов цикл (сертификат — перестановка вершин).
    \item \(P \subseteq NP\). Вопрос \(P = NP\) открыт (скорее всего, нет).
\end{itemize}
\subsection{NP-полнота}
Задача \(L\) называется NP-полной, если:
\begin{enumerate}
    \item \(L \in NP\);
    \item Любая задача из NP сводится к \(L\) за полиномиальное время.
\end{enumerate}
Примеры NP-полных задач из теории графов:
\begin{itemize}
    \item Гамильтонов цикл;
    \item Вершинное покрытие (найти множество вершин минимального размера, покрывающее все рёбра);
    \item Клика (существует ли полный подграф размера \(k\));
    \item Раскраска графа в \(k\) цветов (для \(k \ge 3\)).
\end{itemize}
Эти задачи не имеют известных полиномиальных алгоритмов.

\end{document}